\documentclass[11pt,a4paper]{report}

% Packages
\usepackage{amsmath,amssymb}
\usepackage{bm}
\usepackage{graphicx}
\usepackage{listings}
\usepackage{booktabs}
\usepackage{hyperref}
\usepackage{geometry}
\usepackage{enumitem}
\usepackage{float}
\usepackage{longtable}

\geometry{margin=2.5cm}

\hypersetup{
    colorlinks=true,
    linkcolor=blue,
    citecolor=blue,
    urlcolor=blue
}

\lstset{
    language=Python,
    basicstyle=\ttfamily\small,
    keywordstyle=\color{blue},
    commentstyle=\color{gray},
    stringstyle=\color{red},
    showstringspaces=false,
    breaklines=true,
    frame=single,
    numbers=left,
    numberstyle=\tiny\color{gray}
}

\newcommand{\pd}[2]{\frac{\partial #1}{\partial #2}}
\newcommand{\pdd}[2]{\frac{\partial^2 #1}{\partial #2^2}}
\newcommand{\rvec}{\bm{r}}
\newcommand{\nvec}{\bm{n}}

\title{\textbf{3D Curvilinear Grid Generation}\\[0.5em]
\large Python+C Implementation --- User Manual and Algorithm Reference}
\author{}
\date{\today}

\begin{document}

\maketitle

\begin{abstract}
This document provides a comprehensive reference for the 3D curvilinear grid generation Python+C package. The package implements three grid generation methods --- parabolic, hyperbolic, and elliptic --- for constructing body-fitted 3D meshes suitable for seismic wave numerical simulation. This manual covers the mathematical foundations, algorithmic details of each method, grid quality metrics, post-processing capabilities, and configuration guide.
\end{abstract}

\tableofcontents

%=============================================================================
\chapter{Introduction}
%=============================================================================

\section{Background}

Curvilinear grid generation is essential for finite-difference seismic wave simulation in domains with complex topography. Unlike Cartesian grids, curvilinear grids conform to irregular boundaries, eliminating staircase approximation errors at the free surface.

In 3D, the domain is a hexahedral region bounded by six curved surfaces. The grid generation methods must handle the additional geometric complexity compared to 2D, including six boundary faces, three coordinate directions, and three-dimensional metric tensors.

This package provides Python implementations of three grid generation methods, originally developed in C and documented in a doctoral thesis~[1]:

\begin{itemize}
    \item \textbf{Parabolic method} --- A predictor-corrector scheme that marches from one boundary surface to the opposite, generating grids with controlled clustering near the surface. MPI parallelized in the y-direction. Refer to Chapter~3 of~[1] for the full mathematical derivation.
    \item \textbf{Hyperbolic method} --- A marching method based on orthogonality and volume conservation, suitable for extreme topography (steep cliffs, vertical dams). Serial execution with 3$\times$3 block Thomas algorithm. Refer to Chapter~4 of~[1] for details.
    \item \textbf{Elliptic method} --- An iterative method based on solving elliptic PDEs with source terms (Dirichlet or Higenstock), supporting 3D MPI parallel execution. This method is code-documented only (not covered in~[1]).
\end{itemize}

\section{Scope}

This manual covers:
\begin{itemize}
    \item Mathematical foundations of 3D curvilinear coordinates
    \item Algorithmic principles and implementation details for all three methods
    \item Grid quality metrics (11 metrics)
    \item Post-processing (merging, sampling, stretching, PML)
    \item Configuration and usage guide
\end{itemize}

\section{Computational Coordinates Convention}

Throughout this document, we use the following notation:
\begin{itemize}
    \item Physical coordinates: $(x, y, z)$
    \item Computational coordinates: $(\xi, \eta, \zeta)$, where $\xi$ corresponds to the x-direction, $\eta$ to the y-direction, and $\zeta$ to the z-direction
    \item In the code, $\xi$ is referred to as \texttt{xi}, $\eta$ as \texttt{et}, and $\zeta$ as \texttt{zt}
    \item Grid indices: $i \in [0, N_x-1]$ for $\xi$, $j \in [0, N_y-1]$ for $\eta$, $k \in [0, N_z-1]$ for $\zeta$
    \item Array shape: $(N_z, N_y, N_x)$ in row-major C order, matching NumPy convention
\end{itemize}

\section{Architecture}

The project uses a hybrid Python+C architecture:
\begin{itemize}
    \item \textbf{Python}: Configuration parsing (JSON), memory allocation (NumPy), MPI communication (mpi4py), NetCDF I/O, coordinate permutation
    \item \textbf{C backend} (\texttt{libgrid3d.so}): Compute-intensive kernels loaded via ctypes
\end{itemize}

Array layout follows row-major C order: $\text{index} = k \cdot N_y \cdot N_x + j \cdot N_x + i$.

%=============================================================================
\chapter{Mathematical Foundations}
%=============================================================================

\section{Curvilinear Coordinate Transformation}

The mapping from computational domain to physical domain is:
\begin{equation}
    x = x(\xi, \eta, \zeta), \quad y = y(\xi, \eta, \zeta), \quad z = z(\xi, \eta, \zeta).
\end{equation}

The position vector at each grid point is:
\begin{equation}
    \rvec = [x, y, z]^T.
\end{equation}

\section{Metric Tensors}

In 3D, the covariant metric tensor has 6 independent components:
\begin{equation}
\begin{aligned}
    g_{11} &= \rvec_{,\xi} \cdot \rvec_{,\xi} = x_{,\xi}^2 + y_{,\xi}^2 + z_{,\xi}^2, \\
    g_{22} &= \rvec_{,\eta} \cdot \rvec_{,\eta} = x_{,\eta}^2 + y_{,\eta}^2 + z_{,\eta}^2, \\
    g_{33} &= \rvec_{,\zeta} \cdot \rvec_{,\zeta} = x_{,\zeta}^2 + y_{,\zeta}^2 + z_{,\zeta}^2, \\
    g_{12} &= \rvec_{,\xi} \cdot \rvec_{,\eta} = x_{,\xi} x_{,\eta} + y_{,\xi} y_{,\eta} + z_{,\xi} z_{,\eta}, \\
    g_{13} &= \rvec_{,\xi} \cdot \rvec_{,\zeta} = x_{,\xi} x_{,\zeta} + y_{,\xi} y_{,\zeta} + z_{,\xi} z_{,\zeta}, \\
    g_{23} &= \rvec_{,\eta} \cdot \rvec_{,\zeta} = x_{,\eta} x_{,\zeta} + y_{,\eta} y_{,\zeta} + z_{,\eta} z_{,\zeta}.
\end{aligned}
\end{equation}

The contravariant metric tensor components (used in the SOR update formula) are:
\begin{equation}
\begin{aligned}
    \alpha_1 &= g_{22} g_{33} - g_{23}^2, \\
    \alpha_2 &= g_{11} g_{33} - g_{13}^2, \\
    \alpha_3 &= g_{11} g_{22} - g_{12}^2, \\
    \beta_{12} &= g_{13} g_{23} - g_{12} g_{33}, \\
    \beta_{23} &= g_{12} g_{13} - g_{11} g_{23}, \\
    \beta_{13} &= g_{12} g_{23} - g_{13} g_{22}.
\end{aligned}
\end{equation}

\section{Jacobian Determinant}

The Jacobian determinant represents the volume scaling factor of the 3D coordinate transformation, given by the scalar triple product:
\begin{equation}
    J = \rvec_{,\xi} \cdot (\rvec_{,\eta} \times \rvec_{,\zeta})
      = \begin{vmatrix} x_{,\xi} & x_{,\eta} & x_{,\zeta} \\
                         y_{,\xi} & y_{,\eta} & y_{,\zeta} \\
                         z_{,\xi} & z_{,\eta} & z_{,\zeta} \end{vmatrix}.
\end{equation}

For a valid grid, $J > 0$ must hold everywhere, ensuring no cell folding or inversion.

\section{Finite Difference Approximations}

First derivatives (central differences):
\begin{equation}
\begin{aligned}
    x_{,\xi} &\approx \frac{1}{2}(x_{i+1,j,k} - x_{i-1,j,k}), \\
    x_{,\eta} &\approx \frac{1}{2}(x_{i,j+1,k} - x_{i,j-1,k}), \\
    x_{,\zeta} &\approx \frac{1}{2}(x_{i,j,k+1} - x_{i,j,k-1}).
\end{aligned}
\end{equation}

Second derivatives:
\begin{equation}
\begin{aligned}
    x_{,\xi\xi} &\approx x_{i-1,j,k} - 2x_{i,j,k} + x_{i+1,j,k}, \\
    x_{,\eta\eta} &\approx x_{i,j-1,k} - 2x_{i,j,k} + x_{i,j+1,k}, \\
    x_{,\zeta\zeta} &\approx x_{i,j,k-1} - 2x_{i,j,k} + x_{i,j,k+1}.
\end{aligned}
\end{equation}

Mixed derivatives:
\begin{equation}
\begin{aligned}
    x_{,\xi\eta} &\approx \frac{1}{4}(x_{i+1,j+1,k} + x_{i-1,j-1,k} - x_{i-1,j+1,k} - x_{i+1,j-1,k}), \\
    x_{,\xi\zeta} &\approx \frac{1}{4}(x_{i+1,j,k+1} + x_{i-1,j,k-1} - x_{i-1,j,k+1} - x_{i+1,j,k-1}), \\
    x_{,\eta\zeta} &\approx \frac{1}{4}(x_{i,j+1,k+1} + x_{i,j-1,k-1} - x_{i,j-1,k+1} - x_{i,j+1,k-1}).
\end{aligned}
\end{equation}

%=============================================================================
\chapter{Parabolic Grid Generation}
%=============================================================================

The parabolic method generates grids by marching from one boundary surface toward the opposite, using a predictor-corrector approach. It combines algebraic prediction with elliptic correction at each marching step.

\section{Governing Equation}

The parabolic grid generation equation is a quasi-elliptic PDE in 3D. At each marching layer $k$, the equation enforces smoothness through metric-based correction:
\begin{equation}
    g_{22} g_{33} f_{,\xi\xi} + g_{11} g_{33} f_{,\eta\eta} + g_{11} g_{22} f_{,\zeta\zeta} - 2(g_{12} g_{33} f_{,\xi\eta} + g_{13} g_{22} f_{,\xi\zeta} + g_{23} g_{11} f_{,\eta\zeta}) = 0,
    \label{eq:para-gov}
\end{equation}
where $f$ represents either the $x$, $y$, or $z$ coordinate, and the metric coefficients are computed from the current grid.

\section{Predictor Step: Algebraic Prediction}

The predictor step generates an initial guess for the next grid layer using a combination of orthogonal projection and linear extrapolation.

\subsection{Normal Vector Computation}

In 3D, the unit normal vector to the current marching surface is computed via the cross product of the two tangent vectors:
\begin{equation}
    \nvec = \frac{\rvec_{,\xi} \times \rvec_{,\eta}}{|\rvec_{,\xi} \times \rvec_{,\eta}|},
\end{equation}
where the components are:
\begin{equation}
\begin{aligned}
    v_x &= y_{,\xi} z_{,\eta} - z_{,\xi} y_{,\eta}, \\
    v_y &= z_{,\xi} x_{,\eta} - x_{,\xi} z_{,\eta}, \\
    v_z &= x_{,\xi} y_{,\eta} - y_{,\xi} x_{,\eta}.
\end{aligned}
\end{equation}

For \texttt{t2b = 1} (marching from top to bottom), the normal direction is reversed: $\nvec \leftarrow -\nvec$.

\subsection{Prediction Points}

Two candidate prediction points are computed from the known layer $k-1$:

\textbf{Orthogonal prediction} (normal to current surface):
\begin{equation}
    \rvec_o = \rvec_{i,j}^{k-1} + \varepsilon_c |\bm{R}| \nvec,
\end{equation}
where $\bm{R} = \rvec_{i,j}^{N_z-1} - \rvec_{i,j}^{k-1}$ is the vector from the current point to the target boundary.

\textbf{Linear prediction} (along displacement direction):
\begin{equation}
    \rvec_s = \rvec_{i,j}^{k-1} + \varepsilon_c \bm{R}.
\end{equation}

\textbf{Final prediction} (weighted combination):
\begin{equation}
    \rvec_{i,j}^{k+1} = \varepsilon_s \rvec_o + (1 - \varepsilon_s) \rvec_s,
    \label{eq:para-predict}
\end{equation}
where the switching factor is:
\begin{equation}
    \varepsilon_s = e^{-a \zeta}, \quad \zeta = \frac{k}{N_z - 1}.
\end{equation}

Here, $a$ is the clustering coefficient (parameter \texttt{coef} in the configuration). A smaller $a$ causes $\varepsilon_s$ to decay more slowly, resulting in better grid orthogonality near the free surface, but potentially reducing grid smoothness and even causing grid crossing.

\subsection{Step Coefficients}

The coefficients $\varepsilon_c$ and $\varepsilon_i$ are computed using normalized arc-length distances. Let the cumulative arc length be:
\begin{equation}
    d_0 = 0, \quad d_k = d_{k-1} + h_{k-1} \quad (k = 1, 2, \ldots, N_z-1),
\end{equation}
where $h_k$ is the step length from the step file. Then:
\begin{equation}
\begin{aligned}
    \varepsilon_c &= \frac{d_{k+1} - d_{k-1}}{d_{N_z-1} - d_{k-1}}, \\
    \varepsilon_i &= \frac{d_k - d_{k-1}}{d_{k+1} - d_{k-1}}.
\end{aligned}
\end{equation}

The layer $k$ point is then interpolated between layers $k-1$ and $k+1$:
\begin{equation}
    \rvec_{i,j}^{k} = \rvec_{i,j}^{k-1} + \varepsilon_i (\rvec_{i,j}^{k+1} - \rvec_{i,j}^{k-1}).
\end{equation}

\section{Corrector Step: Elliptic Solve}

After prediction, the corrector step solves the discretized elliptic equation~\eqref{eq:para-gov} using the Thomas algorithm for the resulting tridiagonal system.

\subsection{Discretization}

Discretizing Eq.~\eqref{eq:para-gov} at each interior point $(i, j, k)$ for each marching layer $k$: the known values are $f_{i,j}^{k-1}$ (previous layer) and $f_{i,j}^{k+1}$ (predicted layer). In the $\xi$-direction, this yields a tridiagonal system:
\begin{equation}
    a_i f_{i-1,j}^{k} + b_i f_{i,j}^{k} + c_i f_{i+1,j}^{k} = D_i,
\end{equation}
where the coefficients involve the metric terms computed from the predicted grid.

\subsection{Thomas Algorithm}

The tridiagonal system is solved for $x$, $y$, and $z$ simultaneously using the Thomas algorithm (TDMA), which runs in $O(N)$ time:

\begin{enumerate}
    \item \textbf{Forward elimination}: For $i = 2, 3, \ldots, N-1$:
    \begin{equation}
    \begin{aligned}
        c'_i &= c_i / (b_i - a_i c'_{i-1}), \\
        D'_i &= (D_i - a_i D'_{i-1}) / (b_i - a_i c'_{i-1}),
    \end{aligned}
    \end{equation}
    with $c'_1 = c_1/b_1$ and $D'_1 = D_1/b_1$.

    \item \textbf{Back substitution}: For $i = N-1, N-2, \ldots, 1$:
    \begin{equation}
        f_i = D'_i - c'_i f_{i+1}.
    \end{equation}
\end{enumerate}

The same Thomas solve is then applied in the $\eta$-direction (y-direction) to complete the correction.

\section{Boundary Conditions}

At the lateral boundaries, geometric symmetry is applied:
\begin{equation}
\begin{aligned}
    \rvec_{1,j}^{k} &= 2\rvec_{2,j}^{k} - \rvec_{3,j}^{k}, \\
    \rvec_{N_x-1,j}^{k} &= 2\rvec_{N_x-2,j}^{k} - \rvec_{N_x-3,j}^{k}, \\
    \rvec_{i,1}^{k} &= 2\rvec_{i,2}^{k} - \rvec_{i,3}^{k}, \\
    \rvec_{i,N_y-1}^{k} &= 2\rvec_{i,N_y-2}^{k} - \rvec_{i,N_y-3}^{k}.
\end{aligned}
\end{equation}

\section{Marching Direction}

The parameter \texttt{t2b} controls the marching direction:
\begin{itemize}
    \item \texttt{t2b = 1}: March from top boundary ($k = N_z-1$) to bottom boundary ($k = 0$)
    \item \texttt{t2b = 0}: March from bottom boundary ($k = 0$) to top boundary ($k = N_z-1$)
\end{itemize}

The parameter \texttt{direction} specifies which physical direction corresponds to the marching direction:
\begin{itemize}
    \item \texttt{"z"}: Default, march in the z-direction
    \item \texttt{"x"}: Permute coordinates so that x becomes the marching direction
    \item \texttt{"y"}: Permute coordinates so that y becomes the marching direction
\end{itemize}

\section{MPI Decomposition}

The parabolic method uses a 1D Cartesian MPI topology decomposing only the y-direction:
\begin{itemize}
    \item \textbf{Topology}: 1D Cartesian, $N_{\text{proc}}$ processes along $\eta$
    \item \textbf{Ghost cells}: 1 layer on each side in the y-direction
    \item \textbf{Communication}: Ghost coordinates are exchanged after the predict step and after the update step at each marching layer
\end{itemize}

%=============================================================================
\chapter{Hyperbolic Grid Generation}
%=============================================================================

The hyperbolic method generates grids by marching from a single boundary surface, enforcing orthogonality and volume conservation at each step. It is particularly suitable for extreme topography where the parabolic method may fail.

\section{Governing Equations}

The hyperbolic grid generation is based on two conditions in 3D:

\textbf{Orthogonality conditions} (two equations):
\begin{equation}
    \rvec_{,\xi} \cdot \rvec_{,\zeta} = 0, \quad \rvec_{,\eta} \cdot \rvec_{,\zeta} = 0.
    \label{eq:hyper-orth}
\end{equation}

\textbf{Volume conservation condition}:
\begin{equation}
    \rvec_{,\xi} \times \rvec_{,\eta} \cdot \rvec_{,\zeta} = \Delta V_k,
    \label{eq:hyper-vol}
\end{equation}
where $\Delta V_k$ is the specified cell volume at layer $k$.

These three equations (two orthogonality + one volume) determine the three unknowns $(x, y, z)$ at each new layer.

\section{Linearization and Matrix Form}

Linearizing Eqs.~\eqref{eq:hyper-orth}--\eqref{eq:hyper-vol} about the known layer $k-1$, the system can be written in matrix form:
\begin{equation}
    \bm{A}_{k-1} \rvec_{k,\xi} + \bm{B}_{k-1} \rvec_{k,\eta} + \bm{C}_{k-1} \rvec_{k,\zeta} = \bm{f}_k,
\end{equation}
where $\bm{A}$, $\bm{B}$, $\bm{C}$ are 3$\times$3 coefficient matrices computed from the metric terms of layer $k-1$, and $\bm{f}_k$ contains the volume constraint.

\section{Block Tridiagonal System}

Writing the system in terms of the increment $\Delta \rvec_k = \rvec_k - \rvec_{k-1}$ and applying central differences for $\xi$- and $\eta$-derivatives yields two block tridiagonal systems with 3$\times$3 blocks:

\begin{enumerate}
    \item \textbf{$\eta$-direction solve}: Solve for intermediate variable $\bm{g}$ along the $\eta$-direction:
    \begin{equation}
        \bm{a}_j \bm{g}_{j-1} + \bm{b}_j \bm{g}_j + \bm{c}_j \bm{g}_{j+1} = \bm{d}_j,
    \end{equation}
    where $\bm{a}_j, \bm{b}_j, \bm{c}_j$ are 3$\times$3 blocks.

    \item \textbf{$\xi$-direction solve}: Using the result from the $\eta$-solve, solve for the coordinate increment $\Delta \rvec$ along the $\xi$-direction:
    \begin{equation}
        \bm{a}_i \Delta \rvec_{i-1} + \bm{b}_i \Delta \rvec_i + \bm{c}_i \Delta \rvec_{i+1} = \bm{g}_i.
    \end{equation}
\end{enumerate}

Each system is solved using the block Thomas algorithm, which generalizes the scalar Thomas algorithm to 3$\times$3 block matrices.

\section{Block Thomas Algorithm}

The block Thomas algorithm for a system with $n$ unknowns and 3$\times$3 blocks:

\begin{enumerate}
    \item \textbf{Forward elimination}: For $i = 2, 3, \ldots, n$:
    \begin{equation}
    \begin{aligned}
        \bm{c}'_i &= (\bm{b}_i - \bm{a}_i \bm{c}'_{i-1})^{-1} \bm{c}_i, \\
        \bm{d}'_i &= (\bm{b}_i - \bm{a}_i \bm{c}'_{i-1})^{-1} (\bm{d}_i - \bm{a}_i \bm{d}'_{i-1}),
    \end{aligned}
    \end{equation}
    where the inverse of a 3$\times$3 matrix is computed analytically.

    \item \textbf{Back substitution}: For $i = n-1, n-2, \ldots, 1$:
    \begin{equation}
        \bm{x}_i = \bm{d}'_i - \bm{c}'_i \bm{x}_{i+1}.
    \end{equation}
\end{enumerate}

\section{Artificial Dissipation}

To ensure stability and smoothness, artificial dissipation terms are added. The dissipation coefficient adapts to local grid geometry:
\begin{equation}
    \varepsilon_{e\xi} = \varepsilon_c \cdot N_\xi \cdot S \cdot \hat{d}_\xi \cdot a_\xi,
\end{equation}
where $\varepsilon_c$ is the base coefficient (parameter \texttt{coef} in config), and the four factors are:

\subsection{Aspect Ratio Factor}
\begin{equation}
    N_\xi = \sqrt{\frac{x_{,\zeta}^2 + y_{,\zeta}^2 + z_{,\zeta}^2}{x_{,\xi}^2 + y_{,\xi}^2 + z_{,\xi}^2}}.
\end{equation}

\subsection{Depth Factor}
\begin{equation}
    S = \sqrt{\frac{k-1}{N_z - 1}}.
\end{equation}

\subsection{Clustering Detection Factor}
\begin{equation}
    \hat{d}_\xi = \max[d_\xi^{2/S}, 0.1],
\end{equation}
\begin{equation}
    d_\xi = \frac{|\rvec_{i+1,j}^{k-2} - \rvec_{i,j}^{k-2}| + |\rvec_{i,j}^{k-2} - \rvec_{i-1,j}^{k-2}|}
                  {|\rvec_{i+1,j}^{k-1} - \rvec_{i,j}^{k-1}| + |\rvec_{i,j}^{k-1} - \rvec_{i-1,j}^{k-1}|}.
\end{equation}

\subsection{Angle Factor}

The angle factor detects concave corners where additional dissipation is needed. Define unit direction vectors:
\begin{equation}
\begin{aligned}
    \rvec_\xi^+ &= \frac{\rvec_{i+1,j}^{k-1} - \rvec_{i,j}^{k-1}}{|\rvec_{i+1,j}^{k-1} - \rvec_{i,j}^{k-1}|}, \\
    \rvec_\xi^- &= \frac{\rvec_{i-1,j}^{k-1} - \rvec_{i,j}^{k-1}}{|\rvec_{i-1,j}^{k-1} - \rvec_{i,j}^{k-1}|}.
\end{aligned}
\end{equation}

Compute the angle $\theta$ at the grid point:
\begin{equation}
    \beta = \text{atan2}(-(\rvec_\xi^+ \times \rvec_\xi^-), (\rvec_\xi^+ \cdot \rvec_\xi^-)),
\end{equation}
\begin{equation}
    \theta = \begin{cases} \beta + 2\pi & \text{if } \beta < 0, \\ \beta & \text{if } \beta \ge 0. \end{cases}
\end{equation}

The angle factor is:
\begin{equation}
    a_\xi = \begin{cases} \dfrac{1}{1 - \cos^2(\theta/2)} & \text{if } 0 < \theta < \pi, \\ 1 & \text{if } \pi \le \theta \le 2\pi. \end{cases}
\end{equation}

The same computation applies to the $\eta$-direction, yielding $\varepsilon_{e\eta}$.

\section{Boundary Conditions}

Two types of lateral boundary conditions are supported:

\begin{itemize}
    \item \textbf{Floating} (\texttt{bdry\_x\_itype=1}): Free boundary, grid points float naturally:
    \begin{equation}
        \Delta \rvec_{1,j} = \Delta \rvec_{2,j}, \quad \Delta \rvec_{N_x-1,j} = \Delta \rvec_{N_x-2,j}.
    \end{equation}

    \item \textbf{Cartesian} (\texttt{bdry\_x\_itype=2}): Enforced uniform spacing in the x-direction. The boundary points are placed at evenly spaced x-coordinates, with y and z determined by the floating condition:
    \begin{equation}
        x_{1,j}^k = x_{1,j}^0 + (x_{2,j}^0 - x_{1,j}^0) \cdot \varepsilon_x.
    \end{equation}
\end{itemize}

The same applies for \texttt{bdry\_y\_itype} in the y-direction.

\section{Arc-Length Stretching}

After grid generation, points can be redistributed along each vertical line using arc-length stretching. The cumulative arc length along a vertical grid line is:
\begin{equation}
    s_0 = 0, \quad s_k = s_{k-1} + |\rvec_k - \rvec_{k-1}| \quad (k = 1, 2, \ldots, N_z-1).
\end{equation}

New grid points are placed at specified arc-length positions via linear interpolation.

\section{Volume Specification}

The cell volume at each layer is determined by the step file:
\begin{equation}
    \Delta V_k = \Delta s_{k-1} \cdot dh_k,
\end{equation}
where $\Delta s_{k-1}$ is the surface area element of the current grid surface and $dh_k$ is the step length from the step file.

%=============================================================================
\chapter{Elliptic Grid Generation}
%=============================================================================

The elliptic method generates grids by iteratively solving elliptic PDEs with source terms. Unlike the marching methods, all six boundary surfaces are fixed and the interior grid is determined by the solution of the PDE. The method supports MPI parallel execution for large-scale grids.

\section{Governing Equation}

The elliptic grid generation system with source terms $P$, $Q$, and $R$ is:
\begin{equation}
    g_{11} \rvec_{,\xi\xi} + g_{22} \rvec_{,\eta\eta} + g_{33} \rvec_{,\zeta\zeta}
    + 2g_{12} \rvec_{,\xi\eta} + 2g_{13} \rvec_{,\xi\zeta} + 2g_{23} \rvec_{,\eta\zeta}
    + P\rvec_{,\xi} + Q\rvec_{,\eta} + R\rvec_{,\zeta} = 0.
    \label{eq:ellip-gov}
\end{equation}

The source terms $P$, $Q$, and $R$ control grid spacing and clustering. Without them ($P = Q = R = 0$), the equation reduces to the Laplace equation, producing smooth but potentially non-uniform grids.

\section{Initialization: Transfinite Interpolation (TFI)}

The initial grid is generated using linear transfinite interpolation from the six boundary surfaces. Let $\xi = i/(N_x-1)$, $\eta = j/(N_y-1)$, and $\zeta = k/(N_z-1)$ denote normalized coordinates.

The 3D TFI formula is:
\begin{equation}
    \rvec(\xi, \eta, \zeta) = \bm{U} + \bm{V} + \bm{W} - \bm{UV} - \bm{UW} - \bm{VW} + \bm{UVW},
\end{equation}
where:

$\xi$-direction interpolation (from x-boundaries $\bm{B}_{x1}$, $\bm{B}_{x2}$):
\begin{equation}
    \bm{U}(\xi, \eta, \zeta) = (1-\xi) \bm{B}_{x1}(\eta, \zeta) + \xi \bm{B}_{x2}(\eta, \zeta),
\end{equation}

$\eta$-direction interpolation (from y-boundaries $\bm{B}_{y1}$, $\bm{B}_{y2}$):
\begin{equation}
    \bm{V}(\xi, \eta, \zeta) = (1-\eta) \bm{B}_{y1}(\xi, \zeta) + \eta \bm{B}_{y2}(\xi, \zeta),
\end{equation}

$\zeta$-direction interpolation (from z-boundaries $\bm{B}_{z1}$, $\bm{B}_{z2}$):
\begin{equation}
    \bm{W}(\xi, \eta, \zeta) = (1-\zeta) \bm{B}_{z1}(\xi, \eta) + \zeta \bm{B}_{z2}(\xi, \eta),
\end{equation}

Pairwise bilinear corrections (to avoid double-counting):
\begin{equation}
\begin{split}
    \bm{UV}(\xi, \eta, \zeta) = & (1-\xi)(1-\eta)\bm{B}_{x1}(0,\zeta) + (1-\xi)\eta\bm{B}_{x1}(1,\zeta) \\
    & + \xi(1-\eta)\bm{B}_{x2}(0,\zeta) + \xi\eta\bm{B}_{x2}(1,\zeta),
\end{split}
\end{equation}

and similarly for $\bm{UW}$, $\bm{VW}$.

The trilinear corner correction:
\begin{equation}
\begin{split}
    \bm{UVW}(\xi, \eta, \zeta) = & (1-\xi)(1-\eta)(1-\zeta)\bm{c}_{000} + \xi(1-\eta)(1-\zeta)\bm{c}_{100} \\
    & + (1-\xi)\eta(1-\zeta)\bm{c}_{010} + (1-\xi)(1-\eta)\zeta\bm{c}_{001} \\
    & + \xi\eta(1-\zeta)\bm{c}_{110} + \xi(1-\eta)\zeta\bm{c}_{101} \\
    & + (1-\xi)\eta\zeta\bm{c}_{011} + \xi\eta\zeta\bm{c}_{111},
\end{split}
\end{equation}
where $\bm{c}_{ijk}$ are the 8 corner points shared by the boundary surfaces.

\section{SOR Iteration}

Equation~\eqref{eq:ellip-gov} is solved iteratively using the successive over-relaxation (SOR) method (Gauss-Seidel with relaxation factor $\omega$; currently $\omega = 1.0$).

\subsection{Update Formula}

For each interior point $(k, j, i)$, the new value is computed as:
\begin{equation}
\begin{split}
    \rvec_{k,j,i}^{\text{tmp}} = \frac{1}{2(\alpha_1 + \alpha_2 + \alpha_3)} \bigg[
        & \alpha_1(\rvec_{k,j,i+1} + \rvec_{k,j,i-1}) + \alpha_2(\rvec_{k,j+1,i} + \rvec_{k,j-1,i}) \\
        & + \alpha_3(\rvec_{k+1,j,i} + \rvec_{k-1,j,i}) \\
        & + 2(\beta_{12}\rvec_{,\xi\eta} + \beta_{23}\rvec_{,\eta\zeta} + \beta_{13}\rvec_{,\xi\zeta}) \\
        & + \alpha_1 P \rvec_{,\xi} + \alpha_2 Q \rvec_{,\eta} + \alpha_3 R \rvec_{,\zeta} \bigg],
\end{split}
\end{equation}
\begin{equation}
    \rvec_{k,j,i}^{\text{new}} = \omega \rvec_{k,j,i}^{\text{tmp}} + (1-\omega) \rvec_{k,j,i}^{\text{old}}.
\end{equation}

Gauss-Seidel ordering is used: updated values are immediately used for subsequent points.

\subsection{Convergence Check}

The iteration converges when the maximum normalized residual falls below \texttt{iter\_err}:
\begin{equation}
    \max_i \frac{|\Delta \rvec_i|}{h_i} < \epsilon,
\end{equation}
where $\Delta \rvec = \rvec^{\text{new}} - \rvec^{\text{old}}$, $h_i$ is the local grid spacing, and $\epsilon$ is the convergence tolerance.

The global maximum is computed via \texttt{MPI\_Allreduce} with \texttt{MPI\_MAX}.

\section{Source Term Computation}

Two methods are available for computing boundary source terms: Dirichlet and Higenstock.

\subsection{Dirichlet Method}

The Dirichlet method computes boundary source terms from the metric coefficients and normal derivatives at the boundary.

\textbf{Ghost point calculation}: For physical boundaries, ghost points are computed by projecting the interior gradient onto the boundary normal:
\begin{enumerate}
    \item Compute the tangent vectors along the boundary surface.
    \item Compute the outward normal via the cross product of tangent vectors.
    \item Project the interior gradient: $\Delta \rvec_n = (\Delta \rvec \cdot \nvec) \nvec$.
    \item Set ghost point: $\rvec_{\text{ghost}} = \rvec_{\text{bdry}} - \Delta \rvec_n$.
\end{enumerate}

The boundary source terms are then:
\begin{equation}
\begin{aligned}
    P_{\text{bdry}} &= -\frac{1}{g_{11}} (\rvec_{,\xi} \cdot \rvec_{,\xi\xi}) - \frac{1}{g_{22}} (\rvec_{,\xi} \cdot \rvec_{,\eta\eta}) - \frac{1}{g_{33}} (\rvec_{,\xi} \cdot \rvec_{,\zeta\zeta}), \\
    Q_{\text{bdry}} &= -\frac{1}{g_{11}} (\rvec_{,\eta} \cdot \rvec_{,\xi\xi}) - \frac{1}{g_{22}} (\rvec_{,\eta} \cdot \rvec_{,\eta\eta}) - \frac{1}{g_{33}} (\rvec_{,\eta} \cdot \rvec_{,\zeta\zeta}), \\
    R_{\text{bdry}} &= -\frac{1}{g_{11}} (\rvec_{,\zeta} \cdot \rvec_{,\xi\xi}) - \frac{1}{g_{22}} (\rvec_{,\zeta} \cdot \rvec_{,\eta\eta}) - \frac{1}{g_{33}} (\rvec_{,\zeta} \cdot \rvec_{,\zeta\zeta}).
\end{aligned}
\end{equation}

\subsection{Higenstock Method}

The Higenstock method uses angle-based and distance-based control with \texttt{tanh} feedback to enforce orthogonality and spacing at boundaries.

The angle between grid lines at the boundary is:
\begin{equation}
    \theta = \arccos\left( \frac{\rvec_{,\xi} \cdot \rvec_{,\zeta}}{|\rvec_{,\xi}| |\rvec_{,\zeta}|} \right).
\end{equation}

Source terms are updated iteratively with feedback:

\textbf{Orthogonality control}:
\begin{equation}
    Q \leftarrow Q - a_{\text{higen}} \cdot \tanh\left( \frac{\theta_0 - \theta}{\theta_0} \right),
\end{equation}
where $\theta_0 = \pi/2$ (target orthogonality).

\textbf{Distance control}:
\begin{equation}
    P \leftarrow P + a_{\text{diri}} \cdot \tanh\left( \frac{d_{\text{target}} - d_{\text{actual}}}{d_{\text{target}}} \right),
\end{equation}
where $d_{\text{target}}$ is the desired normal distance and $d_{\text{actual}}$ is the current normal distance from the boundary to the first interior layer, computed via the cross product of tangent vectors and dot product projection.

\section{Interior Source Interpolation}

Boundary source terms are interpolated to the interior using linear and exponential decay weighting. For the six boundary faces ($x_1$, $x_2$, $y_1$, $y_2$, $z_1$, $z_2$):

\begin{equation}
\begin{split}
    P = w_0 \big[ &(1-\xi) P_{x1} + \xi P_{x2} \big]
      + w_0 \big[ &e^{-c_2 \eta} P_{y1} + e^{-c_3(1-\eta)} P_{y2} \big] \\
      + w_1 \big[ &e^{-c_4 \zeta} P_{z1} + e^{-c_5(1-\zeta)} P_{z2} \big],
\end{split}
\end{equation}

and similarly for $Q$ and $R$ with appropriate permutations of the exponential and linear weights. Here, $c = [c_0, c_1, c_2, c_3, c_4, c_5]$ are the \texttt{coef} parameters (6 values for 6 boundaries) and $w = [w_0, w_1]$ are the \texttt{weight} parameters.

The exponential decay ensures that source terms from one boundary diminish toward the opposite boundary, preventing unintended grid distortion far from the controlling boundary.

\section{MPI Domain Decomposition}

The elliptic solver uses a 3D Cartesian topology for MPI parallelization:

\begin{enumerate}
    \item \textbf{Domain decomposition}: The global grid is divided into $N_{px} \times N_{py} \times N_{pz}$ subdomains, each assigned to one MPI rank.
    \item \textbf{Ghost cells}: Each subdomain has one layer of ghost cells on each of the 6 sides, exchanged with neighboring ranks after each SOR sweep.
    \item \textbf{Communication}:
    \begin{itemize}
        \item \texttt{MPI\_Isend}/\texttt{MPI\_Irecv} for ghost cell exchange in 6 directions
        \item \texttt{MPI\_Allreduce(..., MPI\_SUM)} for boundary source term aggregation
        \item \texttt{MPI\_Allreduce(..., MPI\_MAX)} for convergence check
        \item \texttt{MPI\_Bcast} for initial configuration broadcast
    \end{itemize}
    \item \textbf{Load balancing}: All ranks perform the same number of SOR iterations; convergence is checked globally.
\end{enumerate}

%=============================================================================
\chapter{Grid Quality Metrics}
%=============================================================================

\section{Orthogonality}

In 3D, there are three orthogonality angles, one for each pair of coordinate directions. Each measures the deviation from $90^\circ$ between intersecting grid lines:
\begin{equation}
    \theta_{\xi\eta} = \arccos\left( \frac{\rvec_{,\xi} \cdot \rvec_{,\eta}}{|\rvec_{,\xi}| |\rvec_{,\eta}|} \right),
    \quad Q_{\xi\eta} = 90 - |\theta_{\xi\eta} - 90|,
\end{equation}
\begin{equation}
    \theta_{\xi\zeta} = \arccos\left( \frac{\rvec_{,\xi} \cdot \rvec_{,\zeta}}{|\rvec_{,\xi}| |\rvec_{,\zeta}|} \right),
    \quad Q_{\xi\zeta} = 90 - |\theta_{\xi\zeta} - 90|,
\end{equation}
\begin{equation}
    \theta_{\eta\zeta} = \arccos\left( \frac{\rvec_{,\eta} \cdot \rvec_{,\zeta}}{|\rvec_{,\eta}| |\rvec_{,\zeta}|} \right),
    \quad Q_{\eta\zeta} = 90 - |\theta_{\eta\zeta} - 90|.
\end{equation}

Perfect orthogonality yields $Q = 90$. Values below $45$ indicate poor grid quality.

\section{Jacobian Determinant}

The Jacobian represents the signed volume of grid cells (scalar triple product):
\begin{equation}
    J = \rvec_{,\xi} \cdot (\rvec_{,\eta} \times \rvec_{,\zeta}).
\end{equation}

Negative values indicate cell folding (grid inversion). Zero values indicate degenerate cells.

\section{Step Sizes}

The physical step lengths in each direction:
\begin{equation}
\begin{aligned}
    \Delta s_\xi &= |\rvec_{i+1,j,k} - \rvec_{i,j,k}|, \\
    \Delta s_\eta &= |\rvec_{i,j+1,k} - \rvec_{i,j,k}|, \\
    \Delta s_\zeta &= |\rvec_{i,j,k+1} - \rvec_{i,j,k}|.
\end{aligned}
\end{equation}

The minimum step size determines the CFL stability limit for seismic wave simulation.

\section{Smoothness}

The smoothness metric measures the ratio of adjacent step lengths:
\begin{equation}
\begin{aligned}
    S_\xi &= \max\left( \frac{|\rvec_{i+1} - \rvec_i|}{|\rvec_i - \rvec_{i-1}|}, \frac{|\rvec_i - \rvec_{i-1}|}{|\rvec_{i+1} - \rvec_i|} \right), \\
    S_\eta &= \max\left( \frac{|\rvec_{j+1} - \rvec_j|}{|\rvec_j - \rvec_{j-1}|}, \frac{|\rvec_j - \rvec_{j-1}|}{|\rvec_{j+1} - \rvec_j|} \right), \\
    S_\zeta &= \max\left( \frac{|\rvec_{k+1} - \rvec_k|}{|\rvec_k - \rvec_{k-1}|}, \frac{|\rvec_k - \rvec_{k-1}|}{|\rvec_{k+1} - \rvec_k|} \right).
\end{aligned}
\end{equation}

$S = 1$ indicates uniform spacing. Large values indicate abrupt changes in grid spacing.

\section{Aspect Ratio}

The aspect ratio measures the ratio of step sizes in different directions at each grid point:
\begin{equation}
    R = \max\left( \frac{\Delta s_{\max}}{\Delta s_{\min}}, \frac{\Delta s_{\min}}{\Delta s_{\max}} \right),
\end{equation}
where $\Delta s_{\max}$ and $\Delta s_{\min}$ are the largest and smallest step sizes among the three directions. Values close to 1 indicate well-proportioned cells; large values indicate highly stretched cells.

\section{Minimum Distance (CFL Stability)}

The function \texttt{cal\_min\_dist} computes the minimum point-to-face distance across all interior grid cells, which determines the maximum stable time step for the finite-difference scheme via the CFL condition.

For each interior point, distances to the 8 surrounding cell faces are computed, and the minimum is recorded.

%=============================================================================
\chapter{Post-Processing}
%=============================================================================

\section{Grid Import and Merging}

The post-processor reads grid coordinates from MPI-partitioned NetCDF files and optionally merges multiple grids along the x, y, or z direction.

\subsection{Import}

Each partition is identified by its \texttt{px\{i\}\_py\{j\}\_pz\{k\}} filename, containing global index and count attributes that specify the position and size of the partition within the global grid.

\subsection{Merging}

When merging along the x-direction, grids are concatenated horizontally with one overlapping column removed at each junction:
\begin{equation}
    N_x^{\text{total}} = \sum_i N_x^{(i)} - (n_{\text{grids}} - 1).
\end{equation}
Similarly for the y and z directions.

\section{Grid Sampling (Upsampling)}

Grid sampling increases the resolution of an existing grid using tri-linear interpolation. Given a sampling factor $(s_x, s_y, s_z)$, the new grid dimensions are:
\begin{equation}
    N_x' = (N_x - 1) s_x + 1, \quad N_y' = (N_y - 1) s_y + 1, \quad N_z' = (N_z - 1) s_z + 1.
\end{equation}

The interpolated coordinates at sub-grid position $(\alpha, \beta, \gamma)$ within cell $(i, j, k)$ are:
\begin{equation}
\begin{split}
    \rvec(\alpha, \beta, \gamma) = & (1-\alpha)(1-\beta)(1-\gamma) \rvec_{i,j,k} + \alpha(1-\beta)(1-\gamma) \rvec_{i+1,j,k} \\
    & + (1-\alpha)\beta(1-\gamma) \rvec_{i,j+1,k} + (1-\alpha)(1-\beta)\gamma \rvec_{i,j,k+1} \\
    & + \alpha\beta(1-\gamma) \rvec_{i+1,j+1,k} + \alpha(1-\beta)\gamma \rvec_{i+1,j,k+1} \\
    & + (1-\alpha)\beta\gamma \rvec_{i,j+1,k+1} + \alpha\beta\gamma \rvec_{i+1,j+1,k+1}.
\end{split}
\end{equation}

\section{Arc-Length Stretching}

Arc-length stretching redistributes grid points along each grid line according to a specified arc-length distribution. In 3D, stretching can be applied independently in the $\xi$, $\eta$, and $\zeta$ directions, enabling non-uniform spacing (e.g., finer near the surface, coarser at depth).

\section{PML Layer Support}

The post-processor supports PML (Perfectly Matched Layer) partitioning with up to 6 directions (x1, x2, y1, y2, z1, z2). The parameter \texttt{pml\_weight\_2x} doubles the computational weight of PML regions when distributing the grid across MPI ranks for seismic simulation.

%=============================================================================
\chapter{Configuration Guide}
%=============================================================================

\section{JSON Format}

All parameters are specified in a JSON configuration file. Keys prefixed with \texttt{\#} are treated as comments and ignored by the parser.

\section{Common Parameters}

\begin{longtable}{p{5.5cm} p{2cm} p{7cm}}
\toprule
\textbf{Key} & \textbf{Type} & \textbf{Description} \\
\midrule
\texttt{number\_of\_grid\_points\_x} & int & Number of grid points in x direction \\
\texttt{number\_of\_grid\_points\_y} & int & Number of grid points in y direction \\
\texttt{number\_of\_grid\_points\_z} & int & Number of grid points in z direction \\
\texttt{direction} & string & Marching direction: \texttt{"x"}, \texttt{"y"}, or \texttt{"z"} \\
\texttt{geometry\_input\_file} & string & Path to boundary geometry file \\
\texttt{step\_input\_file} & string & Path to step length file \\
\texttt{grid\_export\_dir} & string & Output directory \\
\texttt{grid\_check} & 0/1 & Enable grid quality checks \\
\texttt{check\_orth} & 0/1 & Check orthogonality (3 angles) \\
\texttt{check\_jac} & 0/1 & Check Jacobian \\
\texttt{check\_step\_xi} & 0/1 & Check step size in $\xi$ \\
\texttt{check\_step\_et} & 0/1 & Check step size in $\eta$ \\
\texttt{check\_step\_zt} & 0/1 & Check step size in $\zeta$ \\
\texttt{check\_smooth\_xi} & 0/1 & Check smoothness in $\xi$ \\
\texttt{check\_smooth\_et} & 0/1 & Check smoothness in $\eta$ \\
\texttt{check\_smooth\_zt} & 0/1 & Check smoothness in $\zeta$ \\
\texttt{check\_ratio} & 0/1 & Check aspect ratio \\
\bottomrule
\end{longtable}

\section{Parabolic Parameters}

\begin{longtable}{p{5.5cm} p{2cm} p{7cm}}
\toprule
\textbf{Key} & \textbf{Type} & \textbf{Description} \\
\midrule
\texttt{number\_of\_mpiprocs} & int & MPI processes in y-direction \\
\texttt{step\_input\_file} & string & Step length file ($N_z-1$ lines) \\
\texttt{coef} & int & Clustering coefficient $a$ in $\varepsilon_s = e^{-a\zeta}$ \\
\texttt{t2b} & 0/1 & Marching direction: 1=top$\to$bottom \\
\bottomrule
\end{longtable}

\section{Hyperbolic Parameters}

\begin{longtable}{p{5.5cm} p{2cm} p{7cm}}
\toprule
\textbf{Key} & \textbf{Type} & \textbf{Description} \\
\midrule
\texttt{step\_input\_file} & string & Step length file ($N_z-1$ lines) \\
\texttt{coef} & int & Base dissipation coefficient $\varepsilon_c$ \\
\texttt{t2b} & 0/1 & Marching direction: 1=top$\to$bottom \\
\texttt{flag\_stretch} & 0/1 & Enable arc-length stretching \\
\texttt{bdry\_x\_itype} & int & x boundary condition type: 1=floating, 2=Cartesian \\
\texttt{bdry\_y\_itype} & int & y boundary condition type: 1=floating, 2=Cartesian \\
\texttt{epsilon\_x} & float & x boundary parameter (for Cartesian type) \\
\texttt{epsilon\_y} & float & y boundary parameter (for Cartesian type) \\
\bottomrule
\end{longtable}

\section{Elliptic Parameters}

\begin{longtable}{p{5.5cm} p{2cm} p{7cm}}
\toprule
\textbf{Key} & \textbf{Type} & \textbf{Description} \\
\midrule
\texttt{number\_of\_mpiprocs\_x} & int & MPI processes in x direction \\
\texttt{number\_of\_mpiprocs\_y} & int & MPI processes in y direction \\
\texttt{number\_of\_mpiprocs\_z} & int & MPI processes in z direction \\
\texttt{grid\_method} & object & Sub-method config (\S\ref{sec:method-config}) \\
\bottomrule
\end{longtable}

\subsection{Method Sub-configuration}
\label{sec:method-config}

Choose one of three methods:

\textbf{Linear TFI only} (\texttt{"linear\_tfi"}):
\begin{itemize}
    \item No additional parameters; generates initial grid via transfinite interpolation only.
\end{itemize}

\textbf{Dirichlet} (\texttt{"elli\_diri"}):
\begin{longtable}{p{5.5cm} p{2cm} p{7cm}}
\toprule
\textbf{Key} & \textbf{Type} & \textbf{Description} \\
\midrule
\texttt{coef} & [int$\times$6] & Source term coefficients $[c_0, \ldots, c_5]$ for 6 boundaries \\
\texttt{weight} & [float$\times$2] & Interpolation weights $[w_0, w_1]$ \\
\texttt{iter\_err} & float & Convergence tolerance $\epsilon$ \\
\texttt{max\_iter} & int & Maximum SOR iterations \\
\bottomrule
\end{longtable}

\textbf{Higenstock} (\texttt{"elli\_higen"}):
\begin{longtable}{p{5.5cm} p{2cm} p{7cm}}
\toprule
\textbf{Key} & \textbf{Type} & \textbf{Description} \\
\midrule
\texttt{coef} & [int$\times$6] & Source term coefficients $[c_0, \ldots, c_5]$ for 6 boundaries \\
\texttt{weight} & [float$\times$2] & Interpolation weights $[w_0, w_1]$ \\
\texttt{iter\_err} & float & Convergence tolerance $\epsilon$ \\
\texttt{max\_iter} & int & Maximum SOR iterations \\
\bottomrule
\end{longtable}

\section{Post-Processing Parameters}

\begin{longtable}{p{5.5cm} p{2cm} p{7cm}}
\toprule
\textbf{Key} & \textbf{Type} & \textbf{Description} \\
\midrule
\texttt{input\_grid\_number} & int & Number of input grids \\
\texttt{input\_grid\_info\_0} & object & First grid config (import dir, dimensions, MPI procs, stretch) \\
\texttt{merge\_direction} & string & Merge direction: \texttt{"x"}, \texttt{"y"}, or \texttt{"z"} \\
\texttt{number\_of\_mpiprocs\_out} & [int, int, int] & Output MPI partitioning $[n_x, n_y, n_z]$ \\
\texttt{flag\_sample} & 0/1 & Enable tri-linear sampling \\
\texttt{sample\_factor} & [int, int, int] & Upsampling factors $[s_x, s_y, s_z]$ \\
\texttt{pml\_weight\_2x} & 0/1 & Double PML computation weight \\
\texttt{number\_of\_pml\_x1} & int & PML layers at x1 boundary \\
\texttt{number\_of\_pml\_x2} & int & PML layers at x2 boundary \\
\texttt{number\_of\_pml\_y1} & int & PML layers at y1 boundary \\
\texttt{number\_of\_pml\_y2} & int & PML layers at y2 boundary \\
\texttt{number\_of\_pml\_z1} & int & PML layers at z1 boundary \\
\texttt{number\_of\_pml\_z2} & int & PML layers at z2 boundary \\
\texttt{grid\_export\_dir} & string & Output directory \\
\bottomrule
\end{longtable}

%=============================================================================
\chapter{Installation and Usage}
%=============================================================================

\section{Dependencies}

\begin{itemize}
    \item Python $\geq$ 3.8
    \item NumPy
    \item Numba
    \item netCDF4
    \item mpi4py (parabolic and elliptic methods)
    \item GCC (for building C backend, C99 support)
\end{itemize}

\section{Building the C Backend}

\begin{lstlisting}[language=bash]
cd src_c && bash build.sh
\end{lstlisting}

This compiles \texttt{libgrid3d.so}, which is loaded via Python ctypes at runtime.

\section{Running Grid Generation}

\subsection{Parabolic (MPI)}
\begin{lstlisting}[language=bash]
# Prepare boundary data
cd parabolic/prep_grid_bdry
python creat_bdry_z.py
python creat_step_norm.py

# Generate grid
cd ../generate_grid
bash start.sh
# Or directly:
mpiexec -np 2 python para.py --config-file config.json --verbose 10
\end{lstlisting}

\subsection{Hyperbolic (serial)}
\begin{lstlisting}[language=bash]
# Prepare boundary data
cd hyperbolic/prep_grid_bdry
python creat_bdry_z.py
python creat_step_norm.py

# Generate grid
cd ../generate_grid
bash start.sh
# Or directly:
python hyper.py --config-file config.json --verbose 10
\end{lstlisting}

\subsection{Elliptic (MPI)}
\begin{lstlisting}[language=bash]
# Prepare boundary data (6 faces)
cd elliptic/prep_grid_bdry
python creat_bdry.py

# Generate grid
cd ../generate_grid
bash start.sh
# Or directly:
mpiexec -np 1 python ellip.py --config-file config.json --verbose 10
\end{lstlisting}

\subsection{Post-Processing}
\begin{lstlisting}[language=bash]
cd grid_post_process
bash start.sh
# Or directly:
python post_pro.py --config-file config.json --verbose 10
\end{lstlisting}

\subsection{All Methods}
\begin{lstlisting}[language=bash]
bash start_all.sh
\end{lstlisting}

\section{Visualization}

Visualization scripts are in the \texttt{plotting/} directory. Parameters are set directly in each script (not via command-line arguments).

\subsection{Grid Slice Plot}

Edit parameters in \texttt{draw\_grid.py} and run:
\begin{lstlisting}[language=bash]
cd plotting
python draw_grid.py
\end{lstlisting}

Key parameters:
\begin{itemize}
    \item \texttt{cfs\_file}: Path to the config JSON of the grid to visualize
    \item \texttt{subs}, \texttt{subc}, \texttt{subt}: Start index / count / stride for x, y, z. Setting \texttt{subc[n]=1} selects a plane: n=0 $\to$ YOZ, n=1 $\to$ XOZ, n=2 $\to$ XOY
    \item \texttt{flag\_km}: Convert coordinates to km (0/1)
    \item \texttt{flag\_print}: Save figure to file (0/1)
\end{itemize}

\subsection{Quality Metric Plot}

Edit parameters in \texttt{draw\_quality.py} and run:
\begin{lstlisting}[language=bash]
cd plotting
python draw_quality.py
\end{lstlisting}

Additional parameter:
\begin{itemize}
    \item \texttt{varnm}: Quality variable name (\texttt{orth\_xiet}, \texttt{orth\_xizt}, \texttt{orth\_etzt}, \texttt{jacobi}, \texttt{step\_xi}, \texttt{step\_et}, \texttt{step\_zt}, \texttt{smooth\_xi}, \texttt{smooth\_et}, \texttt{smooth\_zt}, \texttt{ratio})
\end{itemize}

\subsection{Verification}

Verification scripts compare Python+C output against the original C code:
\begin{lstlisting}[language=bash]
cd docs
bash run_verify_all.sh         # All methods
python verify_merge.py          # Merge only
python quick_verify.py          # Quick check
\end{lstlisting}

%=============================================================================
\appendix
\chapter{File Formats}
%=============================================================================

\section{Geometry File (Parabolic/Hyperbolic)}

For parabolic and hyperbolic methods, the geometry file contains two boundary surfaces:
\begin{lstlisting}[language={}]
# nx ny number
51 41
# bz1 coords (bottom surface: ny x nx points)
x y z
...
# bz2 coords (top surface: ny x nx points)
x y z
...
\end{lstlisting}

Each surface has $N_y \times N_x$ lines, with each line containing the $(x, y, z)$ coordinates of one point. Lines starting with \texttt{\#} are comments.

\section{Geometry File (Elliptic)}

For the elliptic method, the geometry file contains six boundary surfaces:
\begin{lstlisting}[language={}]
51
41
21
# x1 face: ny*nz lines of x y z
# x2 face: ny*nz lines of x y z
# y1 face: nx*nz lines of x y z
# y2 face: nx*nz lines of x y z
# z1 face: nx*ny lines of x y z
# z2 face: nx*ny lines of x y z
\end{lstlisting}

The first three lines specify $N_x$, $N_y$, and $N_z$. Then six face blocks follow, each containing the coordinates of the boundary points.

\section{Step File}

The step file specifies the marching step length at each layer:
\begin{lstlisting}[language={}]
# number of step
200
5.0e-02
5.0e-02
...
\end{lstlisting}

The first data line is the number of steps ($N_z - 1$), followed by the step values. Lines starting with \texttt{\#} are comments.

\section{NetCDF Output Format}

Output files follow the naming convention \texttt{\{prefix\}\_px\{i\}\_py\{j\}\_pz\{k\}.nc}, where:
\begin{itemize}
    \item \texttt{prefix}: \texttt{coord} for coordinates, or quality metric name (\texttt{orth\_xiet}, \texttt{jacobi}, etc.)
    \item \texttt{px\{i\}\_py\{j\}\_pz\{k\}}: MPI partition index
\end{itemize}

Each file contains:
\begin{itemize}
    \item Global attributes: \texttt{global\_index\_of\_first\_physical\_points} and \texttt{count\_of\_physical\_points}
    \item Variables: dimensions \texttt{(k, j, i)} with \texttt{float32} type
    \item For coordinate files: variables \texttt{x}, \texttt{y}, and \texttt{z}
    \item For quality files: one variable named after the quality metric
\end{itemize}

\section{Arc-Length File}

The arc-length file specifies the cumulative arc-length distribution for stretching:
\begin{itemize}
    \item $N$ float values (one per line), monotonically increasing
    \item Lines starting with \texttt{\#} are comments
    \item C format: first line is the number of points, followed by $N$ values
\end{itemize}

\chapter*{References}
\addcontentsline{toc}{chapter}{References}

\begin{enumerate}
    \item Doctoral thesis: ``Seismic Wave Numerical Simulation Methods Under Complex Topography Conditions'' --- Chapters 3 (parabolic method) and 4 (hyperbolic method).
\end{enumerate}

\end{document}
